\begin{abstract}
A residual policy learned on top of a physics controller can cut the number of costly measurements a manipulation task needs, but in friction-limited tasks its worst errors happen at breakaway, when friction drops from its static to its kinetic level. We ask when such a learned correction should defer to a friction model. We derive two breakaway shields. For gravitational pivoting, a closed-form grip floor certifies that post-slip travel stays within the target band whenever the static-to-kinetic ratio $r$ is at most an assumed $r_{\max}$. For controlled sliding, the shield bounds the time spent below the arrest force. For pivoting, a certified release from rest exists only if $r<\varphi_{\lim}/\sin\varphi_{\lim}$, so shallow targets are infeasible for realistic $r$. In a pre-registered MuJoCo study with emulated stick--slip, the two tasks split. In sliding, a residual trained under the shield cut overshoot from \N{slide-1.2-B0-over}\% to \N{slide-1.2-B5-over}\% at $r=1.2$ and ended within 5\,mm on 96--98\% of episodes, but kept only part of the unshielded residual's measurement saving; a confirmation set replicated this. In pivoting, the shield found no certified release for exactly the half of the test targets that the feasibility condition predicts. Deferring to the model then left overshoot at \N{pivot-1.1-B5-over}\% at $r=1.1$. Declining prevented overshoot, but only by stopping short. Within the assumed range, certificates were violated on at most 0.5\% of steps, mostly through simulator contact creep. A breakaway shield helps where the physics leaves room for a certified action and, where it does not, makes that limit explicit.
\end{abstract}
